\documentclass[journal]{IEEEtran}
\usepackage{booktabs}
\usepackage{tabularx}
\usepackage{array}
\usepackage[hidelinks]{hyperref}
\title{Supplementary File S6: Integrated Simulator Run}
\author{}
\begin{document}
\maketitle

\noindent This supplement covers only the integrated run reported in
Section~V. The accompanying ZIP contains selected machine-readable records,
source tables, and figure assets; it is not a runtime image.

\section{Run Conditions and Phase Checks}
The run used a deterministic cast-film simulator and tool adapter; no
language-model loop or physical equipment was used. The platform and
simulator were already running, and a shared data directory was reused. The
run identifier was \texttt{20261031034253-1dbo}. The command was
\texttt{node bench/pipeline.mjs --profile integrated --seed 42 --cl-seeds 3}.
The recorded runtime was Node.js 24.19.0 on Windows x64; the recorded source
commit was \texttt{60fa983e}.
All reported measurements refer to this recorded source snapshot; they are not
a fresh validation of the repository's current HEAD.

All 17 phase gates reported PASS. The underlying checks comprised 76 PASS,
one WARN, and zero FAIL outcomes. The warning was the BOPET semantic-card
check in P10; it did not fail the phase gate. Table~\ref{tab:s6-stages}
groups the checks in execution order.

\begin{table}[!t]
\centering
\caption{Grouped check outcomes from the integrated run.}
\label{tab:s6-stages}
\footnotesize
\setlength{\tabcolsep}{3pt}
\begin{tabularx}{\columnwidth}{@{}>{\raggedright\arraybackslash}p{2.65cm}r r r X@{}}
\toprule
Pipeline stages & Pass & Warn & Fail & Step purpose\\
\midrule
P0--P2 & 11 & 0 & 0 & Bootstrap, model reference, and provisioning\\
P3, P4, P4m, P4b--P4f & 30 & 0 & 0 & Acquisition, operation, governance, and records\\
P6--P7 & 7 & 0 & 0 & Seeded objective and multimodal acquisition checks\\
P8, P8b & 16 & 0 & 0 & Second-scene check and bounded-backstop drill\\
P9 & 3 & 0 & 0 & Platform subsystem checks\\
P10 & 9 & 1 & 0 & BOPET line provisioning and scripted task path\\
\midrule
Total & 76 & 1 & 0 & 17 phase gates passed\\
\bottomrule
\end{tabularx}
\end{table}

\section{Acquisition, Write Probes, and Latency}
The primary cast-film scene used five simulated devices across Modbus TCP,
OPC UA, MQTT, HTTP, and Modbus RTU. Five acquisition nodes returned 54 samples.
Modbus RTU was an acquisition
endpoint without a write-control node. Each of the other four protocol paths
received six negative probes; all 24 were rejected by the configured
write-admission checks.
Each of the four writable paths logged one accepted legal write. These are
positive-path checks, not a complete denominator for estimating false-block
rates. The second scene configuration (\texttt{film-line}) was checked within
the same run across five endpoints (Modbus TCP, OPC UA, MQTT, Modbus RTU, and HTTP), returning 12,
12, 5, 11, and 11 samples, respectively (51 total). Its three writable paths
each logged one accepted legal write and rejected three negative probes (9/9
total); no false block or code change was recorded. This is an in-run
configuration check, not an independent deployment or repetition, and the
observations do not estimate a general false-block rate.

\begin{table}[!t]
\centering
\caption{Endpoint checks for the in-run \texttt{film-line} configuration.}
\label{tab:s6-second-scene}
\footnotesize
\setlength{\tabcolsep}{3pt}
\begin{tabularx}{\columnwidth}{@{}>{\raggedright\arraybackslash}p{1.45cm}r r X@{}}
\toprule
Protocol & Samples & Rejected probes & Legal-write check\\
\midrule
Modbus TCP & 12 & 3/3 & One accepted\\
OPC UA & 12 & 3/3 & One accepted\\
MQTT & 5 & 3/3 & One accepted\\
Modbus RTU & 11 & -- & Acquisition only\\
HTTP & 11 & -- & Acquisition only\\
\midrule
Total & 51 & 9/9 & Three accepted checks\\
\bottomrule
\end{tabularx}
\end{table}

\begin{table}[!t]
\centering
\caption{Protocol-level measurements recorded in the integrated run.}
\label{tab:s6-protocol}
\footnotesize
\setlength{\tabcolsep}{3pt}
\begin{tabularx}{\columnwidth}{@{}>{\raggedright\arraybackslash}p{1.35cm}r r X@{}}
\toprule
Protocol & Samples & Rejected probes & Reported $p_{50}/p_{95}$ (ms)\\
\midrule
Modbus TCP & 9 & 6/6 & 47.673/48.348\\
OPC UA & 11 & 6/6 & 21.494/24.175\\
MQTT & 6 & 6/6 & 23.015/94.108\\
HTTP & 13 & 6/6 & 34.414/47.784\\
Modbus RTU & 15 & -- & --\\
\bottomrule
\end{tabularx}
\end{table}

The latency entries are per-protocol summaries stored by the run. The log
describes approximately six writes but does not preserve exact per-protocol
latency denominators or event-level timing observations. The run summary's
31.649\,ms field is the unweighted arithmetic mean of the four per-protocol
$p_{50}$ values, not a pooled median. Latency values are therefore descriptive
and do not support adapter rankings or population-level tail estimates.

\section{Seeded Objective and BOPET Task}
The three fixed-seed objective cases used the simulator's offline reference
$J^*=89.894$. The simulator score is
$J=55J_{\mathrm{thick}}+25J_{\mathrm{quality}}+8J_{\mathrm{energy}}+7J_{\mathrm{throughput}}$.
The thickness term scores deviation from $50\pm2\,\mu\mathrm{m}$ with soft
degradation outside the band; the quality term reflects defect rate, the
energy term screw speed, and the throughput term line speed. The score is
defined only for melt temperature in $195$--$225\,^{\circ}\mathrm{C}$ and
pressure no greater than $22\,\mathrm{MPa}$. The runner reports $J_{end}$ as
the mean of the last two valid
trajectory scores, using the sole valid score when a seed stops before a
second score is available. Table~\ref{tab:s6-seeds} also lists $J_{last}$,
the final per-iteration score, to make the aggregation explicit. Seed 43's
initial observation was inside the configured stopping region, so the run
terminated with zero iterations and writes. This is a zero-update case; it
provides no iterative-convergence evidence. The remaining seed cases used two writes
each.

\begin{table}[!t]
\centering
\caption{Fixed-seed simulator objective values. $J_0$ is the initial score;
$J_{end}$ is the mean of the last two valid trajectory scores, or the sole valid score if the seed stops
before a second score is available; $J_{last}$ is the final per-iteration
score retained in the trajectory.}
\label{tab:s6-seeds}
\footnotesize
\setlength{\tabcolsep}{3pt}
\begin{tabular}{@{}r r r r r r r r@{}}
\toprule
Seed & $J_0$ & $J_{end}$ & $J_{last}$ & $J^*$ & Ratio & Iter. & Writes\\
\midrule
42 & 66.271 & 88.775 & 88.853 & 89.894 & 98.8\% & 2 & 2\\
43 & 86.542 & 86.542 & 86.542 & 89.894 & 96.3\% & 0 & 0\\
44 & 63.981 & 88.384 & 88.540 & 89.894 & 98.3\% & 2 & 2\\
\bottomrule
\end{tabular}
\end{table}

The BOPET trace was 28.05, 27.98, 27.98, and $24.84\,\mu\mathrm{m}$ across the
initial state and three writes; the final reading was within the configured
$25.0\pm0.7\,\mu\mathrm{m}$ band. The task board marked the task COMPLETED,
but its associated parameter-change/optimization record remained open and
unjudged. Treat this as an integration observation, not a completed or causally attributed
optimization result.

\section{Provenance and Interpretation}
The source run lacked a contemporaneous SHA-256 manifest; ZIP hashes verify
delivered files only, not runtime immutability.

\end{document}
